% ------------------------------------------------------------------------
% file `3_thales_similitudes-appliquer-4-solution-body.tex'
%   in folder `xsim/'
%
%     solution of type `appliquer' with id `4'
%
% generated by the `solappliquer' environment of the
%   `xsim' package v0.21 (2022/02/12)
% from source `3_thales_similitudes' on 2023/06/09 on line 155
% ------------------------------------------------------------------------
    \begin{tasks}
        \task $ABC \sim DEF$ ?
        $
            \begin{aligned}[t]
                \dfrac{\abs{AC}}{\abs{DF}}
                 & = \dfrac{2}{4,2} \\
                 & = \dfrac{20}{42} \\
                 & = \dfrac{10}{21} \\
            \end{aligned}
            \hfil
            \begin{aligned}[t]
                \dfrac{\abs{AB}}{\abs{DE}}
                 & = \dfrac{3}{6,3} \\
                 & = \dfrac{30}{63} \\
                 & = \dfrac{10}{21} \\
            \end{aligned}
            \hfil
            \begin{aligned}[t]
                \dfrac{\abs{BC}}{\abs{EF}}
                 & = \dfrac{4}{8,4} \\
                 & = \dfrac{40}{84} \\
                 & = \dfrac{10}{21}
            \end{aligned}
        $\bigbreak

        Par le cas CCC, $ABC \sim DEF$. Donc $\abs{\widehat{D}} = \abs{\widehat{A}} = \ang{104}$.
        \task
        $
            \begin{aligned}[t]
                CD \sslash AB \Rightarrow
                \dfrac{\abs{AE}}{\abs{ED}} & = \dfrac{\abs{EB}}{\abs{EC}}
                \qquad\text{Par le théorème de Thalès}                    \\
                \dfrac{x}{9}               & = \dfrac{3}{5}               \\
                x                          & = 9\cdot\dfrac{3}{5}         \\
                x                          & = \dfrac{18}{5}              \\
                x                          & = 3,6
            \end{aligned}
        $
    \end{tasks}
